\section{Open Problems and Future Directions}
\label{sec:open}

Despite a decade of rapid progress, several fundamental challenges remain open:

\subsection*{O1. Unified Framework for Complex Network Types}

Existing SE variants have been developed independently for directed graphs,
hypergraphs, multi-relational graphs, and dynamic networks.
A pressing open problem is the development of a \emph{unified SE framework}
that handles these structures under a single formalism.
This would require extending the encoding-tree formulation to support
non-standard volume and cut measures, and establishing formal connections
between the specialized variants.

\subsection*{O2. Scalable SE Minimization}

The dominant heuristic SE algorithms scale as $O(n^2)$, limiting
applicability to graphs with hundreds of thousands of nodes.
Although SSSE achieves nearly linear complexity under a sampling approximation,
theoretical guarantees on the approximation quality remain weak.
Developing provably good approximation algorithms for SE minimization, analogous
to those for modularity and the map equation, is an important open direction.

\subsection*{O3. Differentiable SE for Deep Learning Integration}

Current differentiable SE formulations (LSENet, DeSE, Glass-SE) are restricted
to 2-level (flat) trees.
Extending differentiable SE to arbitrary-depth hierarchical trees---while
maintaining well-defined gradients and computational tractability---is an open
problem.
This would enable end-to-end learning of hierarchical representations in GNNs
and Transformers.

\subsection*{O4. SE and Large Language Models}

The recent success of LLMs raises the question of whether SE can be used to
analyse or improve the latent graph structures induced by LLM attention
patterns and knowledge graphs.
Preliminary connections have been drawn via SE-based knowledge
quantification~\cite{wang2023quantifyingKnowledge} and graph-augmented
reasoning~\cite{li2023graphMeetsLLM}, but a principled SE-based theory of LLM
structure is absent.

\subsection*{O5. Theoretical Connections to Other Complexity Measures}

While~\cite{liu2022similaritygap} established a bound between SE and VNE, the
exact relationship between SE and other information-theoretic graph measures
(e.g., Kolmogorov complexity, Fisher information, persistent homology) is not
well understood.
A deeper theoretical investigation would both strengthen SE's foundations and
enable new algorithmic insights.

\subsection*{O6. SE in Partially Observable Environments}

SE-based RL methods (SISA, SIRD, SI2E) have been developed for fully
observable settings.
Extending SE to partial observability---where the agent's state graph is
uncertain or must be inferred---is an important frontier for real-world
applicability.

\subsection*{O7. Benchmark Standardization}

A major bottleneck for the field is the absence of a standardized evaluation
protocol.
Different papers use different datasets, train--test splits, and evaluation
metrics, making direct comparison difficult.
We call for the community to adopt a standard benchmark suite (as partially
initiated in Section~\ref{sec:benchmark}) with open, reproducible
implementations.

\subsection*{O8. SE for Fairness and Privacy}

The use of SE to deliberately obscure community structure~\cite{liu2019rem}
suggests promising applications in \emph{privacy-preserving graph publishing}
and \emph{fairness-aware graph learning}.
A principled theory connecting SE-based obfuscation to differential privacy
guarantees remains an open problem.

